\documentclass[11pt, aspectratio=169]{beamer}
\usepackage[utf8]{inputenc}
\usepackage[T1]{fontenc}
\usepackage{lmodern}
\usepackage[portuguese]{babel}
\usepackage{amsmath}
\usepackage{amssymb}
\usepackage{booktabs}
\usepackage{graphicx}

\usetheme{Madrid}
\usecolortheme{whale}

\title{Controle em Tempo Real}
\subtitle{Controle Ótimo Baseado em Dados via Reinforcement Learning}
\author{Engenharia Matemática \& Controle e Automação}
\date{}

\begin{document}

% Slide 1: Capa
\begin{frame}
    \titlepage
\end{frame}

% Slide 2: Dicionário Controle vs RL
\begin{frame}{A Ponte Matemática: Controle Clássico vs. RL}
    A malha fechada é idêntica em ambas as áreas, mas a formulação passa de analítica para estocástica e baseada em dados:
    
    \vspace{0.3cm}
    \begin{table}
        \centering
        \small
        \renewcommand{\arraystretch}{1.25}
        \begin{tabular}{lll}
            \toprule
            \textbf{Controle Clássico / Moderno} & \textbf{Reinforcement Learning} & \textbf{Notação Matemática} \\
            \midrule
            Planta / Sistema & Ambiente (\textit{Environment}) & $P(s_{t+1} \mid s_t, a_t)$ \\
            Controlador & Agente (\textit{Agent}) & Otimiza a Política $\pi$ \\
            Variáveis de Estado ($x$) & Estado (\textit{State}) & $s_t \in \mathcal{S}$ \\
            Sinal de Controle ($u$) & Ação (\textit{Action}) & $a_t \in \mathcal{A}$ \\
            Função de Custo ($J$) a minimizar & Recompensa ($r$) a maximizar & $r_t = R(s_t, a_t)$ \\
            Lei de Controle ($u = -Kx$) & Política (\textit{Policy}) & $a_t = \pi(s_t)$ ou $\pi(a_t \mid s_t)$ \\
            \bottomrule
        \end{tabular}
    \end{table}
\end{frame}

% Slide 3: NOVO - Política e Recompensa
\begin{frame}{Fundamentos: Política ($\pi$) e Recompensa ($r_t$)}
    \begin{block}{1. Política ($\pi$): A Lei de Controle}
        Mapeia o estado medido $s_t$ no sinal de atuação $a_t$:
        \begin{itemize}
            \item \textbf{Determinística:} $a_t = \mu(s_t)$ (análoga à realimentação clássica $u = -Kx$).
            \item \textbf{Estocástica:} $a_t \sim \pi(\cdot \mid s_t)$ (ex: Gaussiana $\mathcal{N}(\mu_\theta(s_t), \sigma_\theta(s_t))$). Essencial para \textbf{exploração} da dinâmica da planta e robustez contra ruídos.
        \end{itemize}
    \end{block}

    \begin{alertblock}{2. Recompensa ($r_t$): O Custo Instantâneo Invertido}
        Enquanto o controle ótimo minimiza um custo $c_t$, o RL maximiza o sinal escalar $r_t = R(s_t, a_t, s_{t+1})$. No caso quadrático (LQR):
        \[ r_t = -\left( x_t^T Q_c x_t + u_t^T R_c u_t \right) \]
        \textit{Vantagem prática:} Permite embutir restrições físicas de atuadores e termos não-lineares sem exigir solução analítica de Riccati.
    \end{alertblock}
\end{frame}

% Slide 4: NOVO - Funções Valor V(s) e Q(s,a)
\begin{frame}{Funções Valor: $V^\pi(s)$ e $Q^\pi(s,a)$ (O Custo-para-Ir)}
    Avaliam o desempenho acumulado de longo prazo (horizonte descontado por $\gamma \in [0, 1)$):
    
    \vspace{0.2cm}
    \begin{columns}[T]
        \begin{column}{0.48\textwidth}
            \begin{block}{Função Valor de Estado: $V^\pi(s)$}
                Retorno esperado partindo de $s$ e seguindo a lei de controle $\pi$:
                \[ V^\pi(s) = \mathbb{E}_\pi \left[ \sum_{k=0}^{\infty} \gamma^k r_{t+k+1} \,\middle|\, s_t = s \right] \]
                \begin{itemize}
                    \item É o negativo do \textbf{Custo-para-Ir} ($V^*(x) = -J^*(x) = -x^T P x$ no LQR).
                    \item Atua como candidata a \textbf{Função de Lyapunov} ($\Delta V > 0 \iff \Delta J < 0$).
                \end{itemize}
            \end{block}
        \end{column}
        
        \begin{column}{0.48\textwidth}
            \begin{block}{Função Ação-Valor: $Q^\pi(s,a)$}
                Retorno esperado aplicando a ação $a$ em $s$, e seguindo $\pi$ depois:
                \[ Q^\pi(s,a) = \mathbb{E}_\pi \left[ \sum_{k=0}^{\infty} \gamma^k r_{t+k+1} \,\middle|\, s_t = s, a_t = a \right] \]
                \begin{itemize}
                    \item Embute a dinâmica de 1 passo da planta.
                    \item Permite extrair a ação ótima sem conhecer $A$ e $B$: $\pi^*(s) = \arg\max_a Q^*(s,a)$.
                \end{itemize}
            \end{block}
        \end{column}
    \end{columns}
\end{frame}

% Slide 5: Formulação MDP e Bellman
\begin{frame}{Formulação Matemática: MDPs e Equação de Bellman}
    Em vez de equações diferenciais determinísticas $\dot{x} = f(x,u)$, o RL modela a planta como um \textbf{Processo de Decisão de Markov (MDP)}: $(\mathcal{S}, \mathcal{A}, P, R, \gamma)$.
    
    \begin{block}{Objetivo de Otimização (Retorno Esperado)}
        \[ \max_{\pi} \mathbb{E}_{\tau \sim \pi} \left[ \sum_{t=0}^{\infty} \gamma^t R(s_t, a_t) \right], \quad \gamma \in [0, 1) \]
    \end{block}
    
    \begin{alertblock}{Equação de Bellman (Análoga à Hamilton-Jacobi-Bellman / Riccati)}
        \[ Q^{\pi}(s_t, a_t) = \mathbb{E}_{s_{t+1} \sim P} \left[ r_t + \gamma \mathbb{E}_{a_{t+1} \sim \pi} \left[ Q^{\pi}(s_{t+1}, a_{t+1}) \right] \right] \]
    \end{alertblock}
    
    \begin{itemize}
        \item Em \textbf{Deep RL}, a política $\pi_\theta(s)$ e a função valor $Q_\phi(s,a)$ são aproximadas por Redes Neurais Profundas, permitindo controle não-linear em alta dimensionalidade.
    \end{itemize}
\end{frame}

% Slide 6: Arquitetura Actor-Critic
\begin{frame}{Controle Contínuo: Arquitetura Actor-Critic}
    Para sinais de controle analógicos/contínuos ($u \in \mathbb{R}^m$, como torque ou vazão), discretizar o espaço gera explosão combinatória. A solução do estado da arte é a família \textbf{Actor-Critic}:
    
    \vspace{0.3cm}
    \begin{columns}[T]
        \begin{column}{0.48\textwidth}
            \begin{block}{O Ator (Controlador $\pi_\theta$)}
                \begin{itemize}
                    \item Recebe o vetor de estados $s_t$.
                    \item Gera o sinal de controle contínuo $a_t$ diretamente.
                    \item Atualiza seus pesos $\theta$ na direção do gradiente que maximiza a avaliação do Crítico.
                \end{itemize}
            \end{block}
        \end{column}
        
        \begin{column}{0.48\textwidth}
            \begin{block}{O Crítico (Avaliador $Q_\phi$ ou $V_\phi$)}
                \begin{itemize}
                    \item Estima o custo/retorno futuro esperado a partir do estado atual.
                    \item Atualiza seus pesos $\phi$ minimizando o Erro Residual de Bellman (TD-Error).
                    \item Fornece o gradiente de otimização para o Ator.
                \end{itemize}
            \end{block}
        \end{column}
    \end{columns}
\end{frame}

% Slide 7: On-Policy vs Off-Policy
\begin{frame}{On-Policy vs. Off-Policy: O Dilema dos Dados}
    A distinção matemática central reside na distribuição de probabilidade usada para calcular as esperanças matemáticas durante a otimização:
    
    \vspace{0.2cm}
    \begin{block}{1. Abordagem On-Policy (Ex: PPO)}
        \[ \nabla J(\theta) = \mathbb{E}_{\tau \sim \pi_\theta} \left[ \sum_{t=0}^T \nabla_\theta \log \pi_\theta(a_t \mid s_t) A^{\pi_\theta}(s_t, a_t) \right] \]
        \begin{itemize}
            \item Exige trajetórias geradas \textbf{estritamente} pela lei de controle atual $\pi_\theta$.
            \item Após atualizar $\theta$, os dados antigos são descartados. Baixa eficiência de amostra; exige simuladores rápidos (\textit{Sim-to-Real}).
        \end{itemize}
    \end{block}
    
    \begin{block}{2. Abordagem Off-Policy (Ex: SAC, TD3)}
        \[ J(Q) = \mathbb{E}_{(s,a,r,s') \sim \mathcal{D}} \left[ \left( Q(s,a) - \left( r + \gamma \mathbb{E}_{a' \sim \pi}[Q(s',a')] \right) \right)^2 \right] \]
        \begin{itemize}
            \item Recicla dados históricos armazenados em um \textbf{Replay Buffer} ($\mathcal{D}$), gerados por políticas passadas ou controladores clássicos (PID/LQR).
        \end{itemize}
    \end{block}
\end{frame}

% Slide 8: Algoritmos do Estado da Arte
\begin{frame}{Algoritmos do Estado da Arte em Controle Contínuo}
    \begin{itemize}
        \item \textbf{PPO (\textit{Proximal Policy Optimization} - On-Policy):}
        \begin{itemize}
            \item Restringe o tamanho do passo de atualização (\textit{Clipping} / Trust Region).
            \item Garante convergência monotônica e evita instabilidade brusca na malha.
        \end{itemize}
        \vspace{0.25cm}
        \item \textbf{TD3 (\textit{Twin Delayed DDPG} - Off-Policy):}
        \begin{itemize}
            \item Utiliza duas redes Críticas (Duplo Q-Learning) tomando o mínimo entre elas para eliminar a superestimação de valor.
            \item Adiciona ruído de suavização na ação alvo, evitando picos no atuador.
        \end{itemize}
        \vspace{0.25cm}
        \item \textbf{SAC (\textit{Soft Actor-Critic} - Off-Policy):}
        \begin{itemize}
            \item Padrão-ouro atual em robótica. Maximiza recompensa + \textbf{entropia} da política: $\sum \mathbb{E}[r_t + \alpha \mathcal{H}(\pi(\cdot \mid s_t))]$.
            \item Gera políticas estocásticas altamente robustas a perturbações e dinâmicas não modeladas.
        \end{itemize}
    \end{itemize}
\end{frame}

% Slide 9: Comparativo Final e Critério de Projeto
\begin{frame}{Síntese Comparativa e Critério de Engenharia}
    \begin{table}
        \centering
        \resizebox{\textwidth}{!}{%
        \renewcommand{\arraystretch}{1.35}
        \begin{tabular}{lllll}
            \toprule
            \textbf{Algoritmo} & \textbf{Paradigma} & \textbf{Política} & \textbf{Eficiência de Dados} & \textbf{Aplicação Típica em Controle} \\
            \midrule
            \textbf{DDPG} & Off-Policy & Determinística & Alta & Base teórica (sofre com superestimação) \\
            \textbf{TD3} & Off-Policy & Determinística & Alta & Plantas que exigem sinal de controle suave \\
            \textbf{SAC} & Off-Policy & Estocástica & Muito Alta & Robótica física, plantas lentas/custosas \\
            \textbf{PPO} & On-Policy & Estocástica & Baixa & Simuladores paralelos rápidos (\textit{Sim-to-Real}) \\
            \bottomrule
        \end{tabular}%
        }
    \end{table}
    
    \vspace{0.25cm}
    \begin{alertblock}{Regra Prática de Projeto}
        \small
        \begin{itemize}
            \item \textbf{Simulação barata e rápida:} Use \textbf{PPO} pela facilidade de sintonia e estabilidade monotônica.
            \item \textbf{Planta real, física ou simulação pesada (CFD):} Use \textbf{SAC} ou \textbf{TD3} para maximizar o reaproveitamento de dados via \textit{Replay Buffer}.
        \end{itemize}
    \end{alertblock}
\end{frame}

\end{document}